% ======================================================================
% Abstract + keywords (sn-jnl \abstract environment)
% ======================================================================
\abstract{Graph-based traffic forecasters typically rely on an adjacency
matrix --- physical or learned --- that is rarely tested against a
graph-free control. We evaluate GCN and T-GCN backbones on Los-loop and
SZ-Taxi under a single protocol spanning physical, identity-adjacency,
learned contemporaneous, and learned multi-lag graphs with four consumption
mechanisms, four horizons, five seeds, and matched-sparsity and capacity
controls. Graph-free models outperform the physical graph and all single
learned contemporaneous graphs on both benchmarks. The only configuration
that improves over the graph-free control consumes \emph{lag-specific}
multi-lag DAGMA graphs \emph{per timestep}: T-GCN-MultiGSL-Mix reduces RMSE
by up to $14.5\%$ relative to the graph-free baseline on Los-loop at the
first horizon ($5/5$ paired seeds) and by up to $43.0\%$ relative to the
physical T-GCN baseline across horizons. Feeding the identical graphs as a
static union is markedly worse, indicating that utility depends critically
on how the graphs are consumed. On SZ-Taxi, where the learned multi-lag
graphs retain only two edges, all learned variants remain within seed
variability of the graph-free control; at a matched $30$-edge budget, random
and correlation graphs do not reproduce the Los-loop gain. Learned multi-lag
structure is therefore a conditional tool --- worthwhile where the fitted
dependency structure is informative and consumed in a temporally aligned
manner --- and is described throughout as statistical dependency, not
causal influence.}

\keywords{Traffic forecasting \and Graph structure learning \and
Graph convolutional networks \and T-GCN \and Multi-lag graphs \and
Graph consumption}
